% 土星（综述）
% license CCBYSA3
% type Wiki

本文根据 CC-BY-SA 协议转载翻译自维基百科\href{https://en.wikipedia.org/wiki/Saturn}{相关文章}。

\begin{figure}[ht]
\centering
\includegraphics[width=8cm]{./figures/13ddb48d3e7de073.png}
\caption{土星及其显著的行星环，由卡西尼轨道器拍摄\(^\text{[a]}\)} \label{fig_Saturn_1}
\end{figure}
土星是距离太阳第六颗行星，也是太阳系中仅次于木星的第二大行星。它是一颗气态巨行星，其平均半径约为地球的 9 倍。它的平均密度仅为地球的八分之一，但质量却超过地球的 95 倍。尽管土星几乎与木星一样大，但其质量还不到木星的三分之一。土星以 9.59 AU（14.34 亿公里）的距离绕太阳运行，轨道周期为 29.45 年。

土星内部被认为由一颗岩石核心组成，其外包裹着一层深厚的金属氢层，其上为液态氢和液态氦的中间层，最外层为气体外壳。土星呈现淡黄色调，是由于其上层大气中的氨晶体所致。金属氢层中的电流被认为会产生土星的行星磁场，该磁场强度弱于地球，但由于土星体积更大，其磁矩是地球的 580 倍。土星磁场的强度约为木星的二十分之一\(^\text{[27]}\)。土星外层大气整体较为平淡、缺乏对比度，但亦可能出现长期存在的结构。土星的风速可达每小时 1,800 公里（1,100 英里）。

这颗行星拥有明亮且庞大的光环系统，主要由冰粒构成，另含少量岩石碎屑与尘埃。至少有 274 颗卫星绕火星运行，其中 63 颗已有正式命名；这一数字不包含光环中的数百颗小卫星体。土卫六（Titan）是土星最大的卫星，也是太阳系中第二大的卫星，其体积大于（但质量小于）水星，并且是太阳系中唯一拥有稠密大气层的卫星\(^\text{[28]}\)。
\subsection{名称与符号}
土星以罗马财富与农业之神命名，该神是木星（Jupiter）之父。其天文符号（♄）可追溯至希腊的《奥克绪林库斯纸草文稿》（Oxyrhynchus Papyri），在其中该符号可见为希腊字母 κ–ρ（kappa–rho）的连写并带有一条水平短线，作为 Κρονος（Cronus，即土星的希腊名称）的缩写\(^\text{[29]}\)。该符号后演化出类似小写希腊字母 eta 的形状，并在 16 世纪顶部添加十字，以使这一原本的异教符号基督化。

罗马人以土星之名将一周的第七天命名为 Saturday，即 Sāturni diēs，“土星之日”\(^\text{[30]}\)。
\subsection{物理特征}
\begin{figure}[ht]
\centering
\includegraphics[width=8cm]{./figures/91b99dbcbbce021d.png}
\caption{土星与地球及地球月球大小的比较} \label{fig_Saturn_2}
\end{figure}
土星是一颗气态巨行星，主要由氢和氦组成。它缺乏明确的固体表面，但行星内部很可能存在一颗固体核心\(^\text{[31]}\)。由于自转的影响，土星呈现扁球体形状——两极被压扁、赤道区域鼓起。其赤道半径比极半径大超过 10\%：分别为 60,268 公里与 54,364 公里（37,449 英里与 33,780 英里）\(^\text{[6]}\)。

太阳系中的其他巨行星——木星、天王星与海王星——其扁率均低于土星。赤道凸起与自转速率的共同作用导致土星赤道处的有效表面重力仅为 8.96 m/s²，相当于其两极重力的 74\%，且低于地球表面重力。然而，赤道处的逃逸速度接近 36 km/s，远高于地球的逃逸速度\(^\text{[32]}\)。

土星是太阳系中唯一密度低于水的行星——约低 30\%\(^\text{[33]}\)。尽管其核心的密度远高于水，但由于其厚重大气层的存在，整颗行星的平均密度仅为 0.69 g/cm³。木星的质量是地球的 318 倍\(^\text{[34]}\)，而土星的质量是地球的 95 倍\(^\text{[6][35][36][37][38][39]}\)。木星与土星合计占据了太阳系全部行星质量的 92\%\(^\text{[40]}\)。
\subsubsection{内部结构}
\begin{figure}[ht]
\centering
\includegraphics[width=8cm]{./figures/0e870b2d45e0c034.png}
\caption{按比例绘制的土星示意图} \label{fig_Saturn_3}
\end{figure}
尽管土星主要由氢和氦组成，但其大部分质量并非处于气态，因为当密度超过 0.01 g/cm³ 时，氢会成为非理想液体，而这一密度在包含土星 99.9\% 质量的半径范围内即可达到。土星内部的温度、压力与密度在向核心方向不断上升，使得在更深层中氢呈金属态\(^\text{[40]}\)。

标准的行星模型表明，土星内部与木星相似，均包含一颗小型岩石核心，外围包裹着氢和氦，并含有少量各种挥发物\(^\text{[41]}\)。对土星形变的分析显示，其中心物质的集中程度显著高于木星，因此在其中心附近含有更多密度高于氢的物质。土星中心区域按质量计算约有 50\% 为氢，而木星约为 67\%\(^\text{[42]}\)。

该核心在组成上与地球类似，但密度更大。结合土星的引力矩测量结果与其内部物理模型，科学家得以对土星核心的质量设定约束。2004 年的估计认为其核心质量必须为地球质量的 9–22 倍\(^\text{[43][44]}\)，相当于直径约 25,000 公里（16,000 英里）\(^\text{[45]}\)。对土星光环的测量则显示其核心可能更加弥散，质量约等于 17 个地球，半径约为土星整体半径的 60\%\(^\text{[46]}\)。其外层包裹着更厚的液态金属氢层，然后是被氦饱和的液态分子氢层，随高度升高逐渐过渡为气体。最外层约有 1,000 公里（620 英里）厚，由气体组成\(^\text{[47][48][49]}\)。

土星内部非常炽热，其核心温度可达 11,700 °C（21,100 °F），并向外辐射的能量是其自太阳接收能量的 2.5 倍。木星的热能由缓慢的引力压缩（开尔文–亥姆霍兹机制）产生；但这一机制单独可能不足以解释土星的热量来源，因为土星质量更小。另一种或额外的机制可能是土星深层内部的氦滴“降雨”过程。当氦滴向密度较低的氢层下降时，摩擦会释放能量，并使土星外层的氦含量减少\(^\text{[50][51]}\)。这些下降的氦滴可能最终形成包围核心的氦层\(^\text{[41]}\)。在土星内部，如同在木星\(^\text{[52]}\)、以及冰巨行星天王星和海王星\(^\text{[53]}\)中一样，可能会出现钻石雨。
\subsubsection{大气层}  
土星外层大气按体积分数包含 96.3\% 的分子氢和 3.25\% 的氦。与太阳中该元素的丰度相比，土星大气中的氦比例显著偏低\(^\text{[41]}\)。比氦更重的元素（金属丰度）的精确含量尚不明确，但通常认为其比例与太阳系形成时期的原始丰度一致。估计这些较重元素的总质量为地球质量的 19–31 倍，其中相当一部分位于土星的核心区域\(^\text{[54]}\)。

在土星大气中已探测到微量的氨、乙炔、乙烷、丙烷、磷化氢和甲烷\(^\text{[55][56][57]}\)。上层云由氨晶体组成，而下层云似乎由硫氢化铵（NH\(_4\)SH）或水构成\(^\text{[58]}\)。太阳的紫外辐射会在上层大气中引发甲烷的光解，导致一系列烃类化学反应，其生成物再通过湍流与扩散向下输送。此光化学循环受土星年度季节周期的调制\(^\text{[57]}\)。卡西尼号曾观测到一系列出现在北纬的云层结构，被昵称为“珍珠串”。这些结构是位于更深云层中的云隙\(^\text{[59]}\)。

\textbf{云层}
\begin{figure}[ht]
\centering
\includegraphics[width=8cm]{./figures/5b480d0d41a65c33.png}
\caption{2011 年，一场全球性的风暴环绕整个行星。该风暴沿着行星一周移动，使得风暴的头部（明亮区域）绕行并经过其尾部。} \label{fig_Saturn_4}
\end{figure}







